% TOP_PROBLEM: {"id": 600008, "title": "Baker's Dozen — Chains of null geodesics", "category_id": 6, "category": {"id": 6, "name": "geometry", "display_name": "Geometry", "description": "Euclidean and non-Euclidean geometry, geometric structures.", "slug": "geometry", "order_index": 6, "created_at": "2026-07-31T15:26:25.671Z"}, "status": "partially_solved", "problem_number": "AMR-005-0008", "set_id": 13}
% FIRST_POSTED: 2026-10-01
% ENTRY_KIND: statement_only
% UnsolvedMath ID 600008; source catalogue code AMR-005-0008
% !TeX program = lualatex
% Definition and sources only; this file is not an LLM solution attempt.
\documentclass[11pt,a4paper]{article}
\usepackage[margin=25mm]{geometry}
\usepackage{fontspec}
\setmainfont{lmroman10-regular.otf}[BoldFont=lmroman10-bold.otf,
 ItalicFont=lmroman10-italic.otf,BoldItalicFont=lmroman10-bolditalic.otf]
\usepackage{amsmath,amssymb,mathtools}
\usepackage{unicode-math}
\setmathfont{latinmodern-math.otf}
\AtBeginDocument{\let\setminus\smallsetminus}
\usepackage{xurl,hyperref}
\hypersetup{colorlinks=true,urlcolor=blue}
\setcounter{secnumdepth}{0}
\setlength{\parindent}{0pt}
\setlength{\parskip}{5pt}
\setlength{\emergencystretch}{3em}
\begin{document}
\section*{Baker's Dozen — Chains of null geodesics}\label{600008}
{\small UnsolvedMath ID 600008 · AMR-005-0008 · Geometry · AMR Open Problem Lists}
\subsection{Definitions and mathematical statement}
For $a,b,c>0$ consider
\[
 E=\{(x,y,z):x^2/a+y^2/b+z^2/c=1\},\qquad
 g=dx^2+dy^2-dz^2.
\]
The restriction of $g$ to the tangent planes of $E$ degenerates on two closed
curves, called tropics, given by
$z=\pm c\sqrt{x^2/a^2+y^2/b^2}$. Between them the induced metric has signature
$(1,1)$. At every point of this equatorial annulus there are two null tangent
lines, characterized by $g(v,v)=0$; their integral curves are the two families
of unparametrized null geodesics. Label these families left and right
consistently around the annulus.

Starting on a tropic, follow one family to the other tropic, switch families,
and repeat. Each tropic-to-tropic arc counts as one step. An $(n,r)$-chain
returns to its initial point and outgoing family after $n$ steps, and has
winding number $r$ about the equator. Winding is the degree of projection
to the angular coordinate in the $(x,y)$-plane; returning to the same tropic
forces $n$ to be even.

For prescribed $n$ and $r$, find necessary and sufficient conditions on
$(a,b,c)$ for such a chain to exist. A criterion must test these parameters,
for example through explicit equations or periods, rather than merely
rephrasing closure of the trajectory.
\subsection{Short English statement}
Find an explicit closure criterion for alternating lightlike geodesics on
a Minkowski ellipsoid, specifying both the number of steps and the winding.
\subsection{Sources}
\begin{itemize}
\item[S1] ULAM.AI and the UnsolvedMath Contributors, supplied problems.json, record 600008 (AMR-005-0008), AMR Open Problem Lists. Catalogue identity and source transcription; CC BY 4.0.
\url{https://huggingface.co/datasets/ulamai/UnsolvedMath}.
\item[S2] Tabachnikov - A Baker's Dozen of Problems (2015), Problem 4. Original source indexed by AMR-005-0008.
\url{https://amj.math.stonybrook.edu/html-articles/Files-2015-2024/14-01/}.
\end{itemize}
\end{document}
